\documentclass{ieeeaccess}

\usepackage{cite}
\usepackage{amsmath,amssymb,amsfonts}
\usepackage{amsthm}
\usepackage{mathrsfs}
\usepackage{algpseudocode}
\usepackage{graphicx}
\usepackage{textcomp}
%\usepackage[table, dvipsnames]{xcolor}
\usepackage{url}
\usepackage{array}
\usepackage{booktabs}
\usepackage{tabularx}
\usepackage{threeparttable}
\usepackage{makecell}
\usepackage{ragged2e}
\usepackage{multirow}
\usepackage{float}
\usepackage{listings}
\usepackage{stfloats}
\usepackage{caption}
%\usepackage[caption=false,font=normalsize,labelfont=sf,textfont=sf]{subfig}

\def\BibTeX{{\rm B\kern-.05em{\sc i\kern-.025em b}\kern-.08em
    T\kern-.1667em\lower.7ex\hbox{E}\kern-.125emX}}

\begin{document}
\history{Date of publication xxxx 00, 0000, date of current version xxxx 00, 0000.}
\doi{10.1109/ACCESS.2017.DOI}

\title{Beyond Closed-Set Accuracy: Constituent-Level Generalization in Compound Motor Fault Diagnosis}


\author{%
  \uppercase{Evin \c{S}ahin Sadik}\authorrefmark{1},
  \uppercase{H{\"u}seyin Tayyer Canseven}\authorrefmark{2}, \IEEEmembership{Member, IEEE}%
}

\address[1]{Department of Electrical and Electronics Engineering, K\"utahya Dumlup{\i}nar University, 43100, K\"utahya, Turkey (e-mail: evin.sahin@dpu.edu.tr)}
\address[2]{Department of Electrical Engineering, Lappeenranta-Lahti University of Technology (LUT University), 53850, Lappeenranta, Finland}

\tfootnote{This work was supported by the Academy of Finland under the Centre of Excellence Programme for the project High-Speed Electromechanical Energy Conversion Systems.}

\markboth
{\c{S}ahin Sad{\i}k \headeretal: Beyond Closed-Set Accuracy: Constituent-Level Generalization in Compound Motor Fault Diagnosis}
{\c{S}ahin Sad{\i}k \headeretal: Beyond Closed-Set Accuracy: Constituent-Level Generalization in Compound Motor Fault Diagnosis}

\corresp{Corresponding author: H{\"u}seyin Tayyer Canseven (e-mail: huseyin.canseven@lut.fi).}

\begin{abstract}
Reported accuracies for data-driven motor fault diagnosis are frequently near
ceiling, yet they rarely measure the ability that industrial deployment
actually requires: recognizing the elementary fault mechanisms of a compound
fault whose exact combination was never observed during model development.
This paper evaluates that ability on a public multimodal dataset of a 2.2-kW
three-phase induction motor containing 24 diagnostic states, including nine
compound faults that pair a mechanical bearing defect with a second
mechanical, electrical, or electromagnetic mechanism, under twelve speed/load
conditions. The problem is formulated at the constituent level: nine
elementary mechanisms are detected independently from physics-guided,
mechanism-specific vibration- and current-domain features, so unseen
combinations remain expressible. A crossed protocol separates three levels of compound-context exposure (recombination, one-sided, and two-sided context zero-shot) from operating-condition novelty, with run-level, severity-aware splits and all calibration confined to training data. Although closed-set 24-class accuracy reaches 0.913, exact recovery of unseen constituent sets remains between 0.083 and 0.157 for the locked linear probe.
Mechanism-resolved controls show that these failures are strongly
heterogeneous. Dynamic eccentricity is almost perfectly discriminable in isolation (AUROC 0.9995; $F_1=0.923$), yet its compound recall falls to 0.017 and its compound $F_1$ to 0.028. Static eccentricity retains substantial discriminative information but shows marked classifier-family sensitivity, whereas winding is already less specific in isolation and remains difficult across classifier families. The primary context-by-condition interaction is positive but becomes unresolved after matching total training-set cardinality.
Cross-partner transfer is also heterogeneous and includes one statistically supported score-orientation reversal. These results show that constituent-level evaluation exposes mechanism-specific generalization failures that closed-set accuracy alone conceals.

\end{abstract}

\begin{keywords}
compound faults, condition monitoring, constituent-level diagnosis, data
leakage, electrical machines, evaluation protocols, fault diagnosis,
induction motors, multi-label classification, zero-shot learning
\end{keywords}

\titlepgskip=-15pt

\maketitle

\section{Introduction}
\label{sec:introduction}
\PARstart{I}{nduction} motors underpin industrial drivetrains, and their
unplanned failures propagate directly into downtime and safety risk.
Data-driven condition monitoring has consequently produced a large literature
whose reported accuracies now cluster near ceiling on standard laboratory
benchmarks. Deployment experience is harder to reconcile with those numbers:
a model commissioned at one operating point degrades at another, and faults
in service rarely arrive one at a time. A bearing defect that develops
alongside a winding short circuit, broken rotor bars, or eccentricity
confronts the monitoring system with a \emph{compound} state whose exact
combination was, in general, never present in the training history.

Conventional closed-set formulations cannot even express this situation: a
previously unseen combination $A{+}B$ constitutes an unseen class. The
compound-fault literature therefore moved toward decoupling formulations that
predict elementary mechanisms rather than combinations, and, more recently,
toward zero-shot settings in which compound states must be recognized from
single-fault training data alone. As Section~\ref{sec:related} reviews,
however, that literature has evaluated almost exclusively \emph{homogeneous}
compounds---bearing--bearing or bearing--gear pairings whose constituents are
all observable in the same vibration channel---and its reported accuracies of roughly 75--87\% are not directly comparable with
a leakage-safe crossed run-level evaluation, because many representative studies use steady operating conditions and sample- or segment-level evaluation protocols.

This paper asks a deliberately harder question on a public dataset that makes
it answerable. The MCC5-THU motor dataset \cite{chen2026} provides 288
recordings of a 2.2-kW three-phase induction motor across 24 diagnostic
states under twelve speed/load conditions, with synchronized triaxial
vibration, three-phase currents, torque, and a key-phase tachometer. Its nine compound faults include bearing--bearing and bearing-plus-non-bearing combinations involving broken rotor bars, static or dynamic eccentricity, and a stator winding short circuit. Their constituent signatures span vibration- and current-domain evidence, providing a heterogeneous multimodal setting for
constituent-level generalization. To our knowledge, neither this
cross-modal compound regime nor this dataset has previously received a
constituent-level evaluation.

The study is designed so that every claim is bounded by an explicit control.
The formulation is constituent-level: nine elementary mechanisms are detected
independently by regularized linear probes operating on curated,
mechanism-specific physical-signature families derived from angle-resampled
order spectra, envelope-order descriptors, and phase-current sequence
measures. A crossed protocol varies compound-context exposure over three
regimes---recombination, one-sided, and two-sided compound-context
zero-shot---and crosses each with seen and unseen operating conditions, under
run-level, severity-aware splits with all preprocessing, calibration, and
threshold selection confined to training data. The principal constituent-level metrics are accompanied by
input-independent constant-predictor references, so that
prevalence-driven scores cannot masquerade as diagnostic skill; oracle
detectability analyses are paired with recording-sensitivity and same-date
controls, so that separability is not silently attributed to fault
mechanisms; and paired run-clustered bootstrap inference, supplemented by composition-level cluster sensitivity analysis for the context-by-condition interaction, distinguishes supported effects from sampling-unit sensitivity.

This paper makes four contributions. First, it introduces a
constituent-level formulation and a crossed
composition--context--condition protocol for compound motor faults whose
constituents span mechanical and electrical modalities, separating novel
composition, compound-context deprivation, and operating-condition novelty
within a single design (Sections~\ref{sec:methods}--\ref{sec:results}).
Second, it develops a prior-aware, mechanism-resolved diagnostic evaluation
that combines input-independent prevalence references, constituent-level
ranking and thresholded performance, isolated-fault learnability controls, and
classifier-family sensitivity. This decomposition helps distinguish empirical patterns consistent with weak constituent specificity, classifier-capacity sensitivity, or degradation that
emerges when an otherwise learnable constituent is evaluated in compound faults. Third, it presents a controlled analysis of constituent transfer, spanning within-pair oracle detectability, cross-partner transfer, physical-signature preservation, and a statistically supported score-orientation reversal, each bounded by recording-sensitivity and same-date controls that reduce the risk of reading recording-related
variation as fault physics.
Finally, sensitivity and robustness analyses over the decision-threshold
objective, single-fault severity merging, and four classifier families show
that absolute decomposition quality is sensitive to operating point and
classifier family, while the unseen-condition recombination advantage over
two-sided withholding remains directionally positive across all four
classifier families.

The principal findings are that closed-set discriminability (accuracy 0.913)
coexists with exact constituent-set recovery of only 0.083--0.157 on unseen
compounds for the locked probe; and that the resulting failures are strongly
mechanism specific rather than uniform. Bearing-inner, bearing-outer, and
broken-bar constituents remain strongly discriminable, whereas dynamic
eccentricity undergoes a particularly severe transition from near-perfect
isolated-fault discrimination (AUROC 0.9995; $F_1=0.923$) to compound recall
of 0.017 and $F_1$ of 0.028. Static eccentricity retains substantial
discriminative information but is strongly classifier-family sensitive,
whereas winding is already less constituent-specific in isolation and remains
poor across nonlinear alternatives. Cross-partner evidence is often
transferable but can additionally undergo severe context-dependent score
reorientation. The primary crossed analysis suggested a larger penalty from
complete compound-context deprivation when the operating condition was also
unseen, but this interaction became unresolved after matching total
training-set cardinality and is therefore treated as a sensitivity-dependent
directional observation rather than a training-support-independent effect.

The remainder of the paper is organized as follows.
Section~\ref{sec:related} reviews related work and positions the study.
Section~\ref{sec:methods} details the constituent representation, feature
construction, detection, protocol, metrics, and statistical analysis.
Section~\ref{sec:results} reports results and Section~\ref{sec:conclusion}
concludes the paper.


\section{Related Work}
\label{sec:related}

Four research threads intersect in this study: compound-fault decoupling, zero-shot recognition of unseen fault combinations, the reliability of evaluation protocols in data-driven condition monitoring, and generalization across operating conditions. We review each and then position the present work (Table~\ref{tab:positioning}).

Early data-driven treatments of compound faults assigned each observed combination its own class, which cannot express combinations absent from training. Decoupling formulations address this by predicting the elementary mechanisms instead. Huang \emph{et al.} proposed a deep decoupling convolutional network whose multi-label classifier recovers the constituents of gearbox compound faults from raw vibration signals \cite{huang2019ddcnn}, and a one-dimensional CNN variant with a multi-label output layer \cite{huang2019multilabel}. Subsequent work refined the decoupling head: capsule networks with dynamic routing and threshold-based label activation \cite{ltssbow2024}, attentional residual networks combined with active learning for measured (rather than injected) bearing compound faults \cite{decoupling_attn2020}, and mixture-of-experts architectures with task-adapted attention for decoupling compounds unseen as combinations \cite{hemtan2024}. Two properties are common among the representative studies considered here: the constituents are \emph{mechanical} (bearing races, balls, gear teeth), so every constituent is observable in the same vibration channel; and several evaluations rely on steady-speed laboratory recordings and sample- or segment-level protocols that do not explicitly test independence across physical recordings.


A stricter setting withholds all compound data and requires composing single-fault knowledge at test time. Xu \emph{et al.} introduced zero-shot compound-fault diagnosis with manually defined semantic attributes learned from single-fault data \cite{eswa2022}, later replacing hand-crafted attributes with autoencoder-learned fault semantics \cite{eswa2023}. Generative formulations synthesize pseudo-samples of unseen compounds from single-fault semantics \cite{zsl2021,kbs2025}; graph embeddings propagate relational structure among base faults \cite{neuro2025}; and physics-informed semantic spaces anchor the attribute vectors in characteristic-frequency knowledge \cite{song2025gzsl}. Reported accuracies on unseen compounds cluster between roughly 75\% and 87\% \cite{eswa2022,eswa2023,song2025gzsl,appac2025}.

These results, however, concern compounds whose constituents are homogeneous: bearing gear combinations in which both mechanisms modulate the same vibration carrier. The MCC5-THU motor dataset instead contains compound compositions that combine
bearing defects with broken-bar, eccentricity, and winding mechanisms. Their
curated signatures span current-domain carrier and imbalance descriptors as
well as, for eccentricity, shaft-order vibration descriptors. The resulting
setting therefore tests heterogeneous multimodal constituent transfer rather
than a purely vibration-domain compound task. To our knowledge, this cross-modal compound regime has not been evaluated in the zero-shot compound-fault literature, and no published constituent-level results exist for this dataset.


A parallel literature shows that evaluation-protocol choices, more than model choices, drive many reported near-ceiling accuracies. Wheat \emph{et al.} quantified the impact of intra-bearing data leakage in vibration diagnosis and advocated part-to-part splits \cite{wheat2024}; recent work formalizes the distinction between segmentation-level and specimen-level leakage and proposes leakage-safe protocols \cite{gama2025,phme2025}; and open benchmark studies made similar observations for deep architectures across public datasets \cite{zhao2020benchmark}. This thread has concentrated on single-fault classification. Its intersection with compound-fault evaluation is, to our knowledge, unexplored --- yet it is precisely where leakage is most subtle: in the dataset used here, severity variants of one physical combination can leak across a composition
boundary, while compound and single-fault recordings may also differ in recording-related conditions, which can confound naive ``compound versus single'' contrasts. Our protocol design (run-level, severity-aware splits; recording-sensitivity and same-date controls; and constant-predictor references for the principal constituent-level metrics) extends leakage-aware evaluation to the compositional setting.


Diagnosis under variable speed and load has produced a large domain-adaptation and domain-generalization literature, surveyed in \cite{han2025}. On the MCC5-THU family specifically, the dataset authors proposed evidential ensemble learning for real-time multimode diagnosis \cite{liu2024tii} and recently questioned how operating conditions should enter learning-based diagnosis at all \cite{rethinking2025}; other groups applied condition-robust architectures to the companion gearbox release \cite{wang2024,phmap2025}. The dominant protocol in this literature holds one condition out of many \cite{han2025,domainshift2025}. The present study instead evaluates
operating-condition novelty jointly with compound-context exposure through the crossed composition--context--condition design.

\begin{table*}[!t]
\caption{Positioning of this study against representative prior work. $\checkmark$~=~explicitly included; --~=~absent, not applicable, or not reported. ``Cross-modal'' means the compound's constituents are observable in different sensing modalities (mechanical $+$ electrical). ``Prior-aware references'' means metrics are reported against label-prevalence (trivial constant-predictor) references.}
\label{tab:positioning}
\centering
\footnotesize
\setlength{\tabcolsep}{4.5pt}
\begin{tabular}{l l l c c c c c c}
\hline
Study & Task focus & Signals & \makecell{Compound\\faults} &
\makecell{Cross-\\modal} & \makecell{Unseen\\combin.} &
\makecell{Condition\\shift} & \makecell{Leakage-safe\\run-level splits} &
\makecell{Prior-aware\\references} \\
\hline
Huang \emph{et al.} \cite{huang2019ddcnn} & multi-label decoupling & vib. & $\checkmark$ & -- & -- & -- & -- & -- \\
LTSS-BoW-CapsNet \cite{ltssbow2024} & capsule decoupling & vib. & $\checkmark$ & -- & -- & -- & -- & -- \\
HeMTAN \cite{hemtan2024} & unseen-compound decoupling & vib. & $\checkmark$ & -- & $\checkmark$ & -- & -- & -- \\
Xu \emph{et al.} \cite{eswa2022} & zero-shot compounds & vib. & $\checkmark$ & -- & $\checkmark$ & -- & -- & -- \\
Xu \emph{et al.} \cite{eswa2023} & learned fault semantics & vib. & $\checkmark$ & -- & $\checkmark$ & -- & -- & -- \\
Song \emph{et al.} \cite{song2025gzsl} & GZSL, physics semantics & vib. & $\checkmark$ & -- & $\checkmark$ & -- & -- & -- \\
Wheat \emph{et al.} \cite{wheat2024} & leakage in evaluation & vib. & -- & -- & -- & -- & $\checkmark$ & -- \\
Han \emph{et al.} \cite{han2025} & multi-condition survey & multi & -- & -- & -- & $\checkmark$ & -- & -- \\
Liu \emph{et al.} \cite{liu2024tii} & multimode diagnosis & vib.+cur. & -- & -- & -- & $\checkmark$ & -- & -- \\
DRMNN \cite{drmnn2020} & multimodal fusion & vib.+cur. & -- & -- & -- & -- & -- & -- \\
\textbf{This work} & \makecell[l]{constituent-level\\generalization} &
\makecell[l]{vib.+cur.\\(+key-phase)} & $\checkmark$ & $\checkmark$ &
$\checkmark$ & $\checkmark$ & $\checkmark$ & $\checkmark$ \\
\hline
\end{tabular}
\end{table*}

That mechanical faults imprint on stator currents and electrical faults on vibration has long motivated multimodal motor diagnosis \cite{mcsa_review2017}. Deep multimodal fusion typically concatenates modalities into one representation or fuses learned features, e.g.\ multilevel information fusion \cite{multilevel2019}, dynamic-routing multimodal networks \cite{drmnn2020}, adaptive CNN fusion of denoised vibration and current \cite{plos2025}, and comparative studies of the two modalities under deep and classical learners \cite{ayankoso2026}. These works report that multimodal fusion can improve conventional single-fault diagnosis. In the present study, constituent detectors instead use predefined mechanism-specific feature families, restricting each detector to physically relevant vibration- or current-domain descriptors and thereby reducing direct competition between unrelated modality features.

Public rotating-machinery datasets underpin this literature: CWRU \cite{smith2015cwru} and Paderborn \cite{lessmeier2016} for bearings, and more recent multi-fault motor and drivetrain releases such as a PMSM stator-fault dataset \cite{jung2023} and an industrial-scale motor--pump dataset \cite{bruinsma2024}. The MCC5-THU family --- a gearbox release \cite{chen2024gearbox} and the motor release used here \cite{chen2026} --- is distinctive in combining single and compound faults, twelve speed/load profiles with transitional segments, and synchronized vibration, three-phase current, torque, and key-phase channels. The motor release is distributed without constituent-level reference
baselines; the present study provides, to our knowledge, its first systematic constituent-level evaluation.


Table~\ref{tab:positioning} contrasts this study with representative prior work along the axes developed above. The distinguishing combination is: compound faults whose constituents live in different modalities; composition, compound-context, and operating-condition generalization varied jointly in one crossed protocol; run-level, severity-aware, leakage-safe splits; the principal constituent-level metrics interpreted against constituent-specific constant-predictor references so that prior-driven scores cannot masquerade as diagnostic skill; and run-clustered bootstrap inference with multiplicity control.






\section{Materials and Methods}
\label{sec:methods}

The study evaluated constituent-level recovery of compound motor faults under novel composition, compound-context, and operating-condition settings.
Rather than treating each compound state as a separate multiclass label, faults were represented by their elementary mechanisms. The analysis combined physics-guided multimodal features, independent constituent detectors, a crossed generalization protocol, transfer and signature-preservation analyses, and statistical evaluation. All preprocessing, normalization, classifier fitting, and threshold calibration were confined to the corresponding training data to prevent information leakage.


\subsection{Dataset and Constituent-Level Problem Formulation}
\label{subsec:problem_formulation}

The study used the public MCC5-THU motor dataset \cite{chen2026}, comprising 288 physical recordings from a 2.2-kW three-phase induction motor across 24 diagnostic states and 12 speed/load operating conditions. The release contains healthy, single-fault, and nine compound-fault compositions, yielding 108 physical compound runs. Synchronized measurements include triaxial vibration, three-phase current, torque, and keyphase signals; the present constituent analysis used vibration, current, and keyphase measurements. Each experimental run $i$ was represented by a binary constituent vector

\begin{equation}
\mathbf{y}_i =
\left[
y_{i,1},
\ldots,
y_{i,K}
\right]^{\mathrm{T}},
\qquad
\mathbf{y}_i \in \{0,1\}^{K},
\label{eq:multilabel_vector}
\end{equation}

with $K=9$ elementary mechanisms: bearing ball, bearing inner race, bearing
outer race, shaft bending, broken bar, dynamic eccentricity, static
eccentricity, voltage unbalance, and winding fault.

The label cardinality distinguished healthy, single-fault, and compound runs,

\begin{equation}
\|\mathbf{y}_i\|_0 =
\begin{cases}
0, & \text{healthy},\\
1, & \text{single fault},\\
\geq 2, & \text{compound fault}.
\end{cases}
\label{eq:compound_definition}
\end{equation}

The constituent composition of run $i$ was defined as

\begin{equation}
\mathcal{C}_i =
\left\{
f_k : y_{i,k}=1
\right\}.
\label{eq:composition_set}
\end{equation}

Unlike closed-set multiclass diagnosis, where an unseen combination $A+B$
forms an unseen class, the constituent formulation seeks to recover $A$ and
$B$ independently. The resulting generalization problem is therefore whether
constituent knowledge learned from single faults and other compound contexts
transfers to a novel composition.


\subsection{Physics-Guided Multimodal Feature Representation}
\label{subsec:physics_features}

Features were derived from three vibration and three phase-current channels
sampled at $12{,}800$~Hz. Signals were segmented into 2-s analysis windows
using a 1-s stride, corresponding to 50\% overlap. Each window was angle
resampled from the measured keyphase pulses at 256 samples per shaft
revolution before order-domain analysis. For physics-feature extraction, at
most 12 windows were retained per recording; when more were available,
approximately evenly spaced windows spanning the recording were selected.
Window-level descriptors were aggregated to one run-level representation by
their median. Although interquartile-range descriptors were
also generated, the curated constituent detectors used only the median
features and excluded the corresponding \texttt{\_\_iqr} terms.

The run-level multimodal representation was

\begin{equation}
\mathbf{x}_i =
\left[
\left(\mathbf{x}^{(v)}_i\right)^{\mathrm{T}},
\left(\mathbf{x}^{(c)}_i\right)^{\mathrm{T}}
\right]^{\mathrm{T}},
\label{eq:multimodal_feature}
\end{equation}

where $\mathbf{x}^{(v)}_i$ and $\mathbf{x}^{(c)}_i$ denote vibration- and
current-derived descriptors.

For an angle-domain signal $x[n]$, a Hann-window periodogram $P_x(o)$ was
computed. Energy around a target shaft order $o_0$ was accumulated within
$\pm0.12$ order and normalized by the total energy over 0.25--60 orders,

\begin{equation}
R_x(o_0)
=
\frac{
\displaystyle\sum_{|o-o_0|\leq0.12} P_x(o)
}{
\displaystyle\sum_{0.25\leq o\leq60} P_x(o)+\epsilon
}.
\label{eq:implemented_bandratio}
\end{equation}

Both local energy and
$\log_{10}[\max(R_x(o_0),\epsilon)]$ were retained.

For the SKF~6205 bearing, the characteristic orders were

\begin{equation}
o_{\mathrm{BPFO}}=3.585,\qquad
o_{\mathrm{BPFI}}=5.415,\qquad
o_{\mathrm{BSF}}=2.357,
\label{eq:bearing_orders}
\end{equation}

together with their second harmonics. Bearing descriptors were extracted from
the Hilbert-envelope vibration order spectrum. Importantly, the Hilbert
envelope was computed \emph{after} angle resampling; therefore, these
descriptors characterize modulation retained in the 256-samples-per-revolution
angle-domain representation rather than a separately band-pass-filtered
broadband demodulation signal. The angle-domain Nyquist limit is 128 shaft
orders, while the normalized spectral ratios used here were restricted to
0.25--60 orders.

Current descriptors were obtained from the individual phase-current spectra
and the Clarke-space-vector magnitude. The dominant electrical carrier
$o_e$ was identified from the maximum summed three-phase spectral energy over
0.5--30 orders, and descriptors were evaluated at

\begin{equation}
o_e,\qquad o_e\pm1,\qquad o_e\pm2,\qquad 2o_e.
\label{eq:current_targets}
\end{equation}

The $o_e\pm1$ and $o_e\pm2$ terms were treated as carrier-relative
broken-bar-sensitive proxies rather than as a fully slip-parameterized
broken-bar model. Because these sidebands may also vary with load, operating
condition, and other electromagnetic modulation phenomena, they are not
interpreted as mechanism-exclusive slip-frequency measurements.

Current imbalance was characterized by the coefficient of variation of the
three phase RMS values and a phase-order-invariant sequence ratio. For the two
sequence candidates $S_1$ and $S_2$ at $o_e$,

\begin{equation}
r_{\mathrm{seq}}
=
\frac{
\min\left(|S_1|,|S_2|\right)
}{
\max\left(|S_1|,|S_2|\right)+\epsilon
},
\label{eq:sequence_ratio}
\end{equation}

which avoids assuming a predefined phase order and therefore represents
sequence imbalance rather than a prespecified negative-sequence magnitude.

Each constituent detector received a fixed mechanism-specific subset of this
feature bank. Bearing-inner, bearing-outer, and bearing-ball detectors used
BPFI/2BPFI, BPFO/2BPFO, and BSF/2BSF envelope-order features, respectively;
shaft bending used shaft-$1\times$ and shaft-$2\times$ vibration features;
broken bar used carrier-sideband proxies; dynamic and static eccentricity used
shaft-order and carrier-sideband features; voltage unbalance used
sequence-ratio and phase-RMS imbalance; and winding used sequence-ratio,
phase-RMS imbalance, and electrical second-harmonic features. Formally,

\begin{equation}
\mathbf{x}^{(k)}_i = \mathbf{M}_k\mathbf{x}_i,
\label{eq:mechanism_feature_map}
\end{equation}

where $\mathbf{M}_k$ is the predefined selector for constituent $k$. The
resulting subsets contained 16 descriptors for each bearing mechanism and
shaft bending, 40 for broken bar, 56 for each eccentricity mechanism, 2 for
voltage unbalance, and 12 for winding. These dimensions arose from the
predefined channel--order combinations rather than post-hoc optimization, and
no outer-test performance was used to alter the feature definitions or
selector sizes.


\subsection{Constituent Detection and Sensitivity Analyses}
\label{subsec:constituent_detection}

A separate binary detector was trained for each constituent mechanism, allowing
multiple faults to be activated simultaneously without a softmax constraint.
The locked primary model used $L_2$-regularized logistic regression
($C=1$) with class balancing and the predefined mechanism-specific feature
subset $\mathbf{x}^{(k)}_i$ for constituent $k$.

All preprocessing was fitted within the outer-training fold: missing values
were imputed using training-set medians and features were standardized using
training-set statistics. Constituent probabilities were converted to binary
decisions as

\begin{equation}
\hat{y}_{i,k}
=
\mathbb{I}\left(p_{i,k}\geq\tau_k\right),
\label{eq:binary_decision}
\end{equation}

where $\tau_k$ is a constituent-specific threshold.

Thresholds were selected exclusively from outer-training data using
\texttt{GroupKFold} grouped by \texttt{condition\_key}, with four inner folds
whenever the available operating-condition groups permitted and fewer folds
only when necessary. Each operating condition was confined to a single inner
fold. Out-of-fold probabilities from the valid inner splits were pooled for
each constituent, and candidate thresholds from 0.05 to 0.95 in increments of
0.01 were ranked primarily by constituent-level $F_1$, with balanced accuracy
and proximity to 0.5 used as deterministic tie breakers. After threshold
selection, the constituent detector was refitted on the complete outer-training
partition, and the selected threshold was fixed for outer-test evaluation.
No outer-test sample contributed to imputation, standardization, classifier
fitting, or threshold selection.

As post-hoc operating-point sensitivity analyses, threshold calibration was
repeated using $F_{0.5}$ and balanced accuracy as alternative objectives while
keeping the outer splits, preprocessing, fitted detectors, grouped inner
validation, and threshold grid unchanged. These alternatives were used only
to characterize the precision--recall trade-off.

Logistic regression was retained as the locked primary probe to emphasize
constituent transferability rather than architecture optimization. Alternative
classifier families and single-fault severity definitions were evaluated
separately as sensitivity analyses.

Because the primary formulation merges high- and low-severity single-fault
examples, a post-hoc analysis tested whether low-severity support affected the
crossed generalization results. Low-severity single-fault runs for bearing
inner, bearing outer, static eccentricity, and winding were removed (48 runs),
reducing the analysis set to 240 runs while leaving all 108 compound
evaluation runs unchanged. The constituent representation, curated features,
logistic-regression model, crossed outer splits, and grouped inner-CV
threshold calibration were otherwise unchanged. Differences from the locked
H/L-merged analysis were assessed with 5000 paired physical-run
cluster-bootstrap resamples.

The crossed protocol was additionally repeated with a random forest, an
RBF-SVM, and a three-hidden-layer MLP as post-hoc robustness references.
The random forest used 100 trees, square-root feature subsampling, a minimum
leaf size of two, and balanced-subsample class weighting. The RBF-SVM used
$C=1$, \texttt{gamma=scale}, and balanced class weighting. The MLP used 64--32--16 hidden units with ReLU activation, Adam optimization,
$L_2=10^{-4}$, an initial learning rate of $10^{-3}$, and a maximum of 400 iterations; class imbalance was handled using balanced sample weights computed from the outer-training labels.

Across classifier families, constituent labels, curated feature sets, outer
splits, training-only preprocessing, grouped inner cross-validation,
threshold selection, and evaluation metrics were held fixed. No alternative
model or hyperparameter configuration was selected using outer-test
performance.

To determine whether poor compound recovery reflected constituent prevalence,
limited classifier capacity, weak isolated-fault detectability, or degradation
specific to compound contexts, two post-hoc diagnostic analyses were
performed. First, constituent-level performance was summarized with equal
weight across the six crossed evaluation scenarios. For the locked logistic regression, the thresholded constituent metrics were
complemented by AUROC computed from the continuous outer-test scores and by
an input-independent always-positive $F_1$ reference for each constituent.
Corresponding $F_1$ values from RF, RBF-SVM, and MLP were used only as
descriptive classifier-capacity sensitivities and not for post-hoc model
selection.

Second, isolated-fault learnability was evaluated using
leave-one-operating-condition-out testing with the same curated feature sets,
logistic-regression specification, training-only preprocessing, and grouped
inner threshold calibration as the locked pipeline. For each constituent,
isolated single-fault runs were contrasted both with healthy runs and with all
other non-compound states. All runs from the held-out operating condition were
removed before preprocessing, model fitting, and threshold calibration.
Compound runs were excluded entirely from this control. The
single-fault-versus-other-non-compound contrast was used for the primary
isolated-learnability comparison because it tests constituent specificity
rather than separation from health alone.
\subsection{Closed-Set Control Experiment}
\label{subsec:closed_set_control}

An auxiliary 24-class closed-set experiment was used to quantify conventional
in-distribution separability before compositional constraints were imposed.
The reference model was a raw-signal multimodal 1-D CNN with separate
three-channel vibration and current encoders. Each encoder contained four
convolutional blocks (32, 64, 128, and 192 channels), adaptive average pooling,
and a 128-dimensional embedding. The two embeddings were concatenated and
passed through a 256-unit fusion layer and a 24-class softmax output.

The network used 2-s windows sampled at $12{,}800$~Hz and five-fold stratified
cross-validation at the physical-run level, so all windows from a run remained
within the same outer fold. Within each outer-training partition, 20\% of the
training runs were reserved for stratified validation. Channel normalization
parameters were estimated from training data only and applied unchanged to
validation and test windows.

Training used cross-entropy loss and AdamW
($\mathrm{lr}=10^{-3}$, weight decay $10^{-4}$, dropout 0.25), with a maximum
of 25 epochs and early-stopping patience of five. The retained checkpoint
maximized validation run-level macro-$F_1$, with validation loss as the tie
breaker. At inference, window-level softmax probabilities were averaged within
each run and the highest mean probability defined the run-level class.

Final performance was computed from pooled out-of-fold predictions so that
each of the 288 runs was tested exactly once. This experiment served only as
an auxiliary closed-set reference and did not influence model or feature
selection in the constituent-generalization analyses.

\subsection{Crossed Composition--Context--Condition Generalization Protocol}
\label{subsec:crossed_protocol}

A crossed outer protocol separated novel composition, compound-context
exposure, and operating-condition shift. Compound structure was represented
as an undirected graph $G=(V,E)$, where
$V=\{f_1,\ldots,f_K\}$ contains the constituent mechanisms and an edge
$(A,B)\in E$ denotes an observed compound. For a target composition
$\mathcal{C}^{*}=\{A,B\}$, the three context regimes correspond to
increasingly restrictive graph withholding (Fig.~\ref{fig:composition_graph}).

Let $d_{\mathrm{cmp}}(f)$ denote the number of distinct compound contexts
containing constituent $f$ in the outer-training set. The regimes were

\begin{equation}
\begin{aligned}
\text{Recombination:}\quad
&d_{\mathrm{cmp}}(A)>0,\;
d_{\mathrm{cmp}}(B)>0,\\
\text{One-sided:}\quad
&d_{\mathrm{cmp}}(A)=0,\;
d_{\mathrm{cmp}}(B)>0,\\
\text{Two-sided:}\quad
&d_{\mathrm{cmp}}(A)=0,\;
d_{\mathrm{cmp}}(B)=0.
\end{aligned}
\label{eq:context_regimes}
\end{equation}

Recombination withheld only the target pair $A+B$ while retaining both
constituents in other compounds. One-sided context zero-shot removed all
compound exposure for one target constituent while retaining its eligible
single-fault examples; both directions were evaluated. Two-sided context
zero-shot removed compound exposure for both target constituents while
retaining their eligible single-fault examples.

\begin{figure*}[t]
\centering
\includegraphics[width=\textwidth]{Fig1_composition_graph.png}
\caption{Graph interpretation of compound-context withholding. Nodes denote
constituent mechanisms and edges observed compound compositions. Recombination
removes only the target edge, one-sided zero-shot removes compound exposure
for one target constituent, and two-sided zero-shot removes it for both.}
\label{fig:composition_graph}
\end{figure*}

Each context regime was crossed with operating-condition exposure. Let
$\mathcal{O}_{\mathrm{train}}$ denote the set of operating conditions
represented in the outer-training data,

\begin{equation}
c_{\mathrm{test}}\in\mathcal{O}_{\mathrm{train}}
\quad\text{(seen)},\qquad
c_{\mathrm{test}}\notin\mathcal{O}_{\mathrm{train}}
\quad\text{(unseen)}.
\label{eq:condition_regimes}
\end{equation}

For unseen-condition evaluation, all runs acquired under
$c_{\mathrm{test}}$ were removed from the outer-training set before
preprocessing, classifier fitting, and threshold calibration. Eligible runs
from the target condition remained available in the seen-condition setting.
The resulting design comprised three context regimes crossed with two
condition regimes, yielding six principal scenarios.

All splits were defined at the physical-run level before model fitting.
Composition identity was derived from the constituent vector rather than
severity-specific class names, preventing severity variants of the same
composition from crossing the withholding boundary. The directional
one-sided evaluation produced 216 prediction rows from 108 physical compound
runs; recombination and two-sided evaluation each produced 108 rows. The two
one-sided predictions from the same run were kept within the same run-level
bootstrap cluster. The recombination-versus-two-sided context-by-condition
interaction was additionally evaluated using composition-level clustering
(Section~\ref{subsec:statistical_analysis}).

For reproducibility, Algorithm~1 summarizes the outer-split construction.
Context and condition withholding always preceded preprocessing and model
fitting.

\noindent\textbf{Algorithm 1. Outer-split construction for the crossed protocol.}
\begin{algorithmic}[1]
\For{each target compound composition $\mathcal{C}^{*}=\{A,B\}$}
    \For{each compound-context regime}
        \If{recombination}
            \State remove the target compound $A+B$ from the training pool
        \ElsIf{one-sided context zero-shot for $A$}
            \State remove all compound runs containing $A$
            \State retain eligible single-fault runs containing $A$
        \ElsIf{one-sided context zero-shot for $B$}
            \State remove all compound runs containing $B$
            \State retain eligible single-fault runs containing $B$
        \ElsIf{two-sided context zero-shot}
            \State remove all compound runs containing either $A$ or $B$
            \State retain eligible single-fault runs containing $A$ or $B$
        \EndIf
        \For{each operating-condition regime}
            \If{unseen condition}
                \State remove all outer-training runs acquired under
                $c_{\mathrm{test}}$
            \EndIf
            \State fit preprocessing using outer-training data
            \State fit constituent detectors
            \State select thresholds by grouped inner cross-validation
            \State evaluate the held-out target compound run(s)
        \EndFor
    \EndFor
\EndFor
\end{algorithmic}

\begin{table*}[t]
\caption{Outer-split cardinalities under the crossed evaluation protocol.
Training-set sizes are reported as minimum/median/maximum numbers of physical
runs across the corresponding outer folds. The one-sided regime is evaluated
in both directional forms, producing two prediction rows per physical compound
run.}
\label{tab:protocol_cardinality}
\centering
\footnotesize
\begin{tabular}{llcccc}
\toprule
Condition &
Context regime &
Outer folds &
Train runs (min/med/max) &
Test runs/fold &
Prediction rows \\
\midrule

Seen & One-sided      & 18  & 228 / 228 / 264 & 12 & 216 \\
Seen & Recombination  & 9   & 276 / 276 / 276 & 12 & 108 \\
Seen & Two-sided      & 9   & 180 / 216 / 216 & 12 & 108 \\

\midrule

Unseen & One-sided      & 216 & 209 / 209 / 242 & 1 & 216 \\
Unseen & Recombination  & 108 & 253 / 253 / 253 & 1 & 108 \\
Unseen & Two-sided      & 108 & 165 / 198 / 198 & 1 & 108 \\

\bottomrule
\end{tabular}
\end{table*}

The withholding regimes also produced different amounts of training support.
Recombination used 276 and 253 training runs per seen- and unseen-condition
fold, respectively, whereas the corresponding median two-sided sizes were
216 and 198. Median positive support for the target constituents decreased
from 71 to 23.5 under seen conditions and from 65 to 22 under unseen
conditions. Thus, context withholding also induces a training-support
difference, which is examined separately in a size-matched sensitivity analysis.

To assess whether the recombination advantage could be explained by total
training-set cardinality alone, recombination training sets were repeatedly
downsampled within operating-condition groups to match the corresponding
two-sided training-set sizes. The original test sets, constituent
representation, curated features, preprocessing, logistic-regression model,
and grouped inner-CV threshold calibration were retained unchanged.

Downsampling was constrained to preserve at least one non-target compound
context for each target constituent, so that the defining recombination
exposure was maintained. Thirty repeated size-matched subsamples were
evaluated. Uncertainty due to both composition sampling and training-set
subsampling was summarized using 5000 composition-level bootstrap resamples
over the nine target compositions, with a size-matched repeat sampled within
each bootstrap replicate. This analysis controls total training-set
cardinality and operating-condition distribution, but does not equate
target-positive support, which is part of the compound-context exposure
difference by design.

\subsection{Evaluation Metrics}
\label{subsec:evaluation_metrics}

Performance was evaluated at the constituent-set level. Constituent recall was

\begin{equation}
R_{\mathrm{const}}
=
\frac{
\sum_{i=1}^{N}\sum_{k=1}^{K} y_{i,k}\hat{y}_{i,k}
}{
\sum_{i=1}^{N}\sum_{k=1}^{K} y_{i,k}
},
\label{eq:constituent_recall}
\end{equation}

and exact-match recovery was

\begin{equation}
\mathrm{EM}
=
\frac{1}{N}
\sum_{i=1}^{N}
\mathbb{I}
\left(
\hat{\mathbf{y}}_i=\mathbf{y}_i
\right).
\label{eq:exact_match}
\end{equation}

False additions per prediction row were quantified as

\begin{equation}
\mathrm{FA/run}
=
\frac{1}{N}
\sum_{i=1}^{N}\sum_{k=1}^{K}
(1-y_{i,k})\hat{y}_{i,k}.
\label{eq:false_additions}
\end{equation}

Recovery of at least one and all true constituents was defined as

\begin{align}
R_{\mathrm{any}}
&=
\frac{1}{N}\sum_{i=1}^{N}
\mathbb{I}
\left(
\sum_{k=1}^{K} y_{i,k}\hat{y}_{i,k}>0
\right),
\label{eq:any_recovered}\\
R_{\mathrm{all}}
&=
\frac{1}{N}\sum_{i=1}^{N}
\mathbb{I}
\left(
\sum_{k=1}^{K} y_{i,k}\hat{y}_{i,k}
=
\sum_{k=1}^{K} y_{i,k}
\right).
\label{eq:all_recovered}
\end{align}

Micro-$F_1$ was computed over all constituent decisions,

\begin{equation}
F_{1,\mathrm{micro}}
=
\frac{2TP}{2TP+FP+FN}.
\label{eq:micro_f1}
\end{equation}

Macro-$F_1$ was additionally reported over only the constituents represented
positively in the corresponding evaluation subset,

\begin{equation}
F_{1,\mathrm{active}}
=
\frac{1}{|\mathcal{K}_{+}|}
\sum_{k\in\mathcal{K}_{+}} F_{1,k},
\label{eq:active_macro_f1}
\end{equation}

where $\mathcal{K}_{+}$ contains constituents with at least one positive test
instance.

\subsection{Within-Pair Detectability, Cross-Partner Transfer, and Signature Preservation}
\label{subsec:transferability}

To distinguish failed constituent recovery from absence of measurable
evidence, constituent $A$ in a target compound $A+B$ was first evaluated by
the within-pair contrast

\begin{equation}
A+B \quad \mathrm{vs.} \quad B,
\label{eq:oracle_contrast}
\end{equation}

with discriminability quantified as
$\mathrm{AUC}^{\mathrm{within}}_{A|B}$.

Cross-partner transfer then tested whether evidence for $A$ learned from
eligible compound contexts other than $A+B$ generalized to the same target
contrast. The target pair was excluded from training, and the resulting
fixed-orientation score was summarized by
$\mathrm{AUC}^{A|B}_{\mathrm{dir}}$.

Because directional AUC below 0.5 may indicate reversed ranking rather than
loss of separability, an orientation-free diagnostic quantity was also
reported,

\begin{equation}
\begin{aligned}
\mathrm{AUC}_{\mathrm{sep}}
&=
\max\left(
\mathrm{AUC}_{\mathrm{dir}},
1-\mathrm{AUC}_{\mathrm{dir}}
\right)\\
&=
0.5+
\left|
\mathrm{AUC}_{\mathrm{dir}}-0.5
\right|.
\end{aligned}
\label{eq:separability_auc}
\end{equation}

$\mathrm{AUC}_{\mathrm{sep}}$ characterizes rank separability only and is not
deployable performance, because post-hoc score inversion would require
target-label information. The fixed-orientation transfer gap was

\begin{equation}
\Delta_{\mathrm{transfer}}
=
\mathrm{AUC}^{\mathrm{within}}_{A|B}
-
\mathrm{AUC}^{A|B}_{\mathrm{dir}},
\label{eq:transfer_gap}
\end{equation}

and therefore reflects degradation of learned score orientation rather than
loss of orientation-free separability.

Recording-related nuisance variation was assessed separately. Acquisition
timestamps were recovered for 287 of 288 runs across seven recording dates.
The only duplicated fault--condition state recorded twice was used as a
nuisance control, and same-date compound pairs sharing one constituent were
compared under matched operating conditions. These analyses reduced
recording-date mismatch but were treated as sensitivity controls rather than
causal estimates of acquisition-session effects.

A complementary feature-level analysis examined how predefined physical
signatures changed when constituent $A$ was embedded in a compound fault.
Only mechanism-specific curated features defined independently of the
cross-partner results were used.

For each curated feature $j$ of constituent $A$, the reference direction was
estimated exclusively from matched healthy and single-fault runs:

\begin{equation}
d_j^{A}
=
\operatorname{sign}^{+}
\left[
\operatorname{median}_{c}
\left(
x_{j,c}^{A}-x_{j,c}^{H}
\right)
\right],
\label{eq:signature_direction}
\end{equation}

where $c$ indexes matched operating conditions, $x_{j,c}^{H}$ and
$x_{j,c}^{A}$ denote the healthy and single-fault feature values,
respectively, and $\operatorname{sign}^{+}(0)=+1$. A robust
healthy-reference scale $s_j^{A}$ was estimated from
$\{x_{j,c}^{H}\}_c$ using $1.4826$ times the median absolute deviation.
If this estimate was degenerate, the interquartile range divided by $1.349$
and then the standard deviation were used sequentially; if all remained
degenerate, the deterministic fallback
$\max(10^{-3}\operatorname{median}_{c}|x_{j,c}^{H}|,10^{-9})$
was used. No compound observation was used to determine $d_j^{A}$ or
$s_j^{A}$. When multiple recordings shared the same original fault label and operating
condition, their numeric physics features were median-aggregated before forming the matched reference.

The directed standardized evidence at operating condition $c$ was

\begin{equation}
\delta_{j,c}^{A}
=
d_j^{A}
\frac{
x_{j,c}^{A}-x_{j,c}^{H}
}{
s_j^{A}
},
\qquad
\delta_{j,c}^{A|B}
=
d_j^{A}
\frac{
x_{j,c}^{A|B}-x_{j,c}^{H}
}{
s_j^{A}
},
\label{eq:directed_standardized_evidence}
\end{equation}

where $x_{j,c}^{A|B}$ denotes the corresponding feature value when
constituent $A$ occurs with partner $B$. Feature preservation was then
defined as

\begin{equation}
\rho_{j,c}^{A|B}
=
\frac{
\delta_{j,c}^{A|B}
}{
\delta_{j,c}^{A}
}.
\label{eq:preservation_ratio}
\end{equation}

Preservation ratios were evaluated only when
$|\delta_{j,c}^{A}|\geq0.10$ to avoid unstable ratios for negligible
single-fault evidence. For the reported sign-reversal and
severe-suppression rates, the analysis was further restricted to
feature--condition observations with positive single-fault evidence
($\delta_{j,c}^{A}>0$). Negative preservation values indicate feature-level
sign reversal, values between zero and one suppression, values near one
preservation, and values above one amplification. Severe suppression was
defined as $\rho_{j,c}^{A|B}<0.5$. The analysis characterizes feature-level
heterogeneity and is not interpreted as a causal decomposition of the
multivariate classifier score.
\subsection{Statistical Analysis}
\label{subsec:statistical_analysis}

Uncertainty of the principal constituent-level metrics was estimated by
nonparametric cluster bootstrap at the physical-run level. In the one-sided
regime, both directional predictions from the same run were retained within
the same bootstrap cluster. Micro-$F_1$ was treated as the principal aggregate
metric, and the recombination-versus-two-sided difference-of-differences as
the principal context-by-condition contrast. Constituent recall and $R_{\mathrm{all}}$ were supportive secondary outcomes.
Because $R_{\mathrm{all}}$ does not penalize false additions, it was interpreted
jointly with exact match and FA/run rather than as a stand-alone measure.
$R_{\mathrm{any}}$ was retained only descriptively because an input-independent
predictor can attain one. Other mechanism-specific, isolated-fault-learnability, threshold,
classifier-family, and transfer analyses were interpreted as descriptive,
sensitivity, or mechanistic analyses and were not used for post-hoc selection
of the primary model or operating point.

Because the same physical compound runs contributed to the compared context
regimes, paired bootstrap contrasts were evaluated separately under seen and
unseen operating conditions:

\begin{align}
\Delta_{\mathrm{R-O}} &= M_{\mathrm{recomb}}-M_{\mathrm{one}},\\
\Delta_{\mathrm{R-T}} &= M_{\mathrm{recomb}}-M_{\mathrm{two}},\\
\Delta_{\mathrm{O-T}} &= M_{\mathrm{one}}-M_{\mathrm{two}},
\end{align}

where $M$ denotes the metric of interest. Context-by-condition interaction was
quantified as

\begin{equation}
\Delta_{\mathrm{int}}
=
\Delta_{\mathrm{unseen}}
-
\Delta_{\mathrm{seen}}.
\end{equation}

Run-level confidence intervals used 5000 paired cluster-bootstrap resamples
and the 2.5th and 97.5th percentiles; an effect was considered resolved when
the interval excluded zero.

Because the 108 compound runs represent only nine distinct compositions, the
recombination-versus-two-sided interaction was additionally subjected to a
composition-level cluster-bootstrap sensitivity analysis. Each of 10,000
replicates sampled the nine compositions with replacement while retaining all
12 operating-condition runs of each sampled composition and applying identical
composition multiplicities to the seen and unseen cells. Sensitivity to
individual compositions was also examined by leave-one-composition-out
analysis.

Closed-set confidence intervals used 5000 stratified run-level bootstrap
resamples within the 24 diagnostic classes.

For cross-partner transfer, score-orientation reversal was tested using

\begin{equation}
T=
\left|
\mathrm{AUC}_{\mathrm{dir}}-0.5
\right|.
\label{eq:polarity_statistic}
\end{equation}

Directional-AUC confidence intervals used 2000 bootstrap resamples clustered
by operating condition. For the orientation test, positive and negative scores
were swapped within matched operating conditions over 5000 permutation
replicates using
$|\mathrm{AUC}_{\mathrm{dir}}-0.5|$ as the test statistic.
Family-wise multiplicity across the 18 directed constituent--partner
relations was controlled by the Holm procedure separately within each
condition-mode family. A
relation was classified as a statistically supported score-orientation
reversal only when its directional AUC and bootstrap interval were below
0.5 and the Holm-adjusted permutation test supported departure from chance
orientation.

\section{Results}
\label{sec:results}

\subsection{Closed-Set Discrimination and the Generalization Gap}
\label{subsec:closed_set_results}

The auxiliary 24-class closed-set classifier achieved an accuracy of 0.9132,
balanced accuracy of 0.9134, and macro-$F_1$ of 0.9119
(Table~\ref{tab:closed_set_results}), indicating strong conventional
in-distribution separability.

\begin{table}[t]
\caption{Closed-set multiclass diagnostic performance with 95\% bootstrap
confidence intervals.}
\label{tab:closed_set_results}
\centering
\begin{tabular}{lcc}
\toprule
Metric & Estimate & 95\% CI \\
\midrule
Accuracy & 0.9132 & [0.8854, 0.9375] \\
Balanced accuracy & 0.9134 & [0.8856, 0.9377] \\
Macro-$F_1$ & 0.9119 & [0.8826, 0.9366] \\
\bottomrule
\end{tabular}
\end{table}

The substantially lower constituent-level performance under the crossed
protocols is therefore consistent with the additional demands of composition,
compound-context, and operating-condition generalization rather than poor
baseline separability alone. Because the closed-set reference uses a different
model architecture from the locked constituent probe, this comparison is
descriptive rather than a model-controlled estimate of task difficulty;
classifier-family robustness analyses below separately assess the contribution
of constituent-model choice.


\subsection{Crossed Composition--Context--Condition Generalization}
\label{subsec:crossed_results}

The crossed protocol was evaluated on all 108 physical compound-fault runs
under recombination, one-sided, and two-sided context regimes, each with seen
and unseen operating conditions. Both directional removals were evaluated in
the one-sided regime. Table~\ref{tab:crossed_results} summarizes the principal
constituent-level results.

\begin{table*}[t]
\caption{Constituent-level performance across context and operating-condition
regimes. EM denotes exact constituent-set recovery, $R_{\mathrm{const}}$
constituent recall, FA/run false additions per prediction row, and
$R_{\mathrm{any}}$ and $R_{\mathrm{all}}$ recovery of at least one and all true constituents, respectively. Boldface denotes the descriptively best value within each operating-condition block and does not imply a statistically resolved difference.}
\label{tab:crossed_results}
\centering
\footnotesize
\begin{tabular}{llccccccc}
\toprule
Condition &
Context regime &
Active Macro-$F_1$ &
Micro-$F_1$ &
EM &
$R_{\mathrm{const}}$ &
$\mathrm{FA/run}$ &
$R_{\mathrm{any}}$ &
$R_{\mathrm{all}}$ \\
\midrule

Seen
& One-sided
& 0.5890
& 0.6581
& 0.1204
& 0.6505
& 0.6528
& 0.9259
& 0.3750 \\

Seen
& Recombination
& \textbf{0.6177}
& \textbf{0.6775}
& 0.1389
& \textbf{0.6759}
& \textbf{0.6389}
& 0.9352
& \textbf{0.4167} \\

Seen
& Two-sided
& 0.5994
& 0.6620
& 0.1111
& 0.6620
& 0.6759
& \textbf{0.9444}
& 0.3796 \\

\midrule

Unseen
& One-sided
& 0.5761
& 0.6452
& 0.1111
& 0.6481
& 0.7222
& 0.9398
& 0.3565 \\

Unseen
& Recombination
& \textbf{0.6177}
& \textbf{0.6865}
& \textbf{0.1574}
& \textbf{0.6944}
& \textbf{0.6574}
& \textbf{0.9444}
& \textbf{0.4444} \\

Unseen
& Two-sided
& 0.5588
& 0.6311
& 0.0833
& 0.6296
& 0.7315
& 0.9352
& 0.3241 \\

\bottomrule
\end{tabular}
\end{table*}

\begin{table*}[t]
\caption{Input-independent constant-predictor references on the 108 compound
runs, quantifying performance obtainable from constituent prevalence alone.}
\label{tab:trivial_baselines}
\centering
\footnotesize
\begin{tabular}{lcccccc}
\toprule
Constant prediction &
Micro-$F_1$ &
EM &
$R_{\mathrm{const}}$ &
FA/run &
$R_{\mathrm{any}}$ &
$R_{\mathrm{all}}$ \\
\midrule
Inner + outer
& 0.5556 & 0.1111 & 0.5556 & 0.8889 & 1.0000 & 0.1111 \\

Inner + outer + broken bar
& 0.5333 & 0.0000 & 0.6667 & 1.6667 & 1.0000 & 0.3333 \\

Inner only
& 0.3704 & 0.0000 & 0.2778 & 0.4444 & 0.5556 & 0.0000 \\
\bottomrule
\end{tabular}
\end{table*}

Recombination yielded the strongest overall performance, with micro-$F_1$ of
0.6775 and 0.6865 and constituent recall of 0.6759 and 0.6944 under seen and
unseen conditions, respectively. Exact constituent-set recovery nevertheless
remained low across all six regimes (0.0833--0.1574), indicating that partial
recovery was substantially easier than exact decomposition.

The prevalence references in Table~\ref{tab:trivial_baselines} show why
partial-recovery metrics cannot be interpreted in isolation. A constant
inner-plus-outer predictor achieves $R_{\mathrm{any}}=1$, while adding broken
bar raises constituent recall to 0.6667 but produces 1.6667 false additions
per run and no exact matches. By comparison, the locked detector achieved micro-$F_1$ of 0.631--0.686 with
FA/run of 0.639--0.732. Under recombination, $R_{\mathrm{all}}$ reached
0.4167 and 0.4444, compared with 0.3333 for the constant
inner-plus-outer-plus-broken-bar reference. That reference, however, incurred
1.6667 false additions per run, compared with 0.6389 and 0.6574 for the
learned detector. Thus, the advantage of the learned detector is a more
favorable joint recovery--overdiagnosis trade-off rather than uniform
superiority on every individual metric. 
\subsection{Mechanism-Specific Failure Modes: Isolated Learnability,
Compound Degradation, and Model Capacity}
\label{subsec:mechanism_failure_modes}
\begin{table*}[t]
\caption{Mechanism-specific diagnostic decomposition. Single-fault AUROC and
$F_1$ are from leave-one-operating-condition-out discrimination of the target
single fault against all other non-compound states. Compound AUROC, recall,
and LR $F_1$ are averaged equally across the six crossed scenarios.
Const.-$F_1$ is the constituent-specific input-independent always-positive
reference. RF, RBF-SVM, and MLP results are descriptive classifier-family
sensitivities and were not used for model selection.}
\label{tab:constituent_failure_modes}
\centering
\footnotesize
\setlength{\tabcolsep}{3pt}
\begin{tabular}{lccccccccc}
\toprule
Constituent &
\makecell{Single\\AUROC} &
\makecell{Single\\$F_1$} &
\makecell{Compound\\AUROC} &
\makecell{Compound\\Recall} &
\makecell{LR\\$F_1$} &
\makecell{Const.\\$F_1$} &
\makecell{RF\\$F_1$} &
\makecell{RBF-SVM\\$F_1$} &
\makecell{MLP\\$F_1$} \\
\midrule

Bearing inner
& 0.999 & 0.957 & 0.965 & 0.896 & 0.940 & 0.714 & 0.947 & 0.839 & 0.936 \\

Bearing outer
& 1.000 & 1.000 & 0.988 & 0.796 & 0.886 & 0.714 & 0.898 & 0.900 & 0.874 \\

Broken bar
& 1.000 & 1.000 & 0.993 & 0.927 & 0.923 & 0.364 & 0.919 & 0.883 & 0.906 \\

Dynamic eccentricity
& 1.000 & 0.923 & 0.590 & 0.017 & 0.028 & 0.364 & 0.385 & 0.054 & 0.085 \\

Static eccentricity
& 0.992 & 0.920 & 0.867 & 0.538 & 0.577 & 0.364 & 0.658 & 0.866 & 0.814 \\

Winding
& 0.823 & 0.400 & 0.398 & 0.229 & 0.205 & 0.364 & 0.019 & 0.163 & 0.115 \\

\bottomrule
\end{tabular}
\end{table*}
Table~\ref{tab:constituent_failure_modes} reveals that poor compound
decomposition does not arise from a single common failure mechanism.
Bearing-inner, bearing-outer, and broken-bar constituents were highly
discriminable both in isolation and in compound evaluation. Their compound
AUROCs remained between 0.965 and 0.993, and their locked $F_1$ values
substantially exceeded the corresponding input-independent references.
Thus, their strong thresholded performance reflects genuine constituent
discrimination rather than prevalence alone.

Dynamic eccentricity showed the clearest compound-context degradation.
It was almost perfectly discriminable from other non-compound states under
leave-one-condition-out evaluation (AUROC=0.9995; $F_1=0.923$), yet its
compound recall fell to 0.017 and its locked compound $F_1$ to 0.028.
Random forest increased compound $F_1$ to 0.385, indicating additional
classifier-capacity sensitivity, but recovery remained modest. The combination
of strong isolated learnability and severe compound degradation is consistent
with a pronounced loss of decision consistency in the compound setting and
shows that the poor compound result cannot be explained by weak
isolated-fault detectability alone.

Static eccentricity exhibited a different failure mode. It was also strongly learnable in isolation (AUROC=0.992; $F_1=0.920$) and retained substantial
ranking information in compound evaluation (AUROC=0.867), despite only
moderate locked $F_1$ (0.577). The RBF-SVM increased descriptive compound
$F_1$ to 0.866 under the same crossed protocol, indicating pronounced
classifier-family sensitivity rather than absence of discriminative
information in the curated representation.

Winding differed from both eccentricity mechanisms. Its isolated-fault specificity was already weaker (AUROC=0.823; $F_1=0.400$), and compound performance deteriorated further ($F_1=0.205$). None of the nonlinear alternatives improved on the locked linear probe, while
its mean compound AUROC was below chance (0.398). This aggregate
constituent-level AUROC is distinct from the directed cross-partner transfer
analysis below and is therefore not interpreted as evidence of a statistically
supported score-orientation reversal. Rather, it indicates poor fixed-score
ranking on average across the crossed scenarios. Together, these results
reveal three distinct empirical failure patterns: severe compound degradation
for dynamic eccentricity, strong classifier-family sensitivity for static
eccentricity, and weaker constituent specificity with persistent compound
difficulty for winding.

\subsection{Paired Effects of Compound-Context Deprivation}
\label{subsec:context_condition_results}

The principal comparison was the recombination-versus-two-sided
context-by-condition interaction in micro-$F_1$, with constituent recall and
$R_{\mathrm{all}}$ as supportive recovery outcomes. Paired run-cluster
bootstrap contrasts are summarized in
Fig.~\ref{fig:paired_protocol_differences}.

Under seen operating conditions, recombination showed only limited advantages.
For constituent recall, the recombination--one-sided difference was 0.0255
(95\% CI: $[0.0000,\,0.0509]$), whereas the recombination--two-sided
difference was 0.0139 (95\% CI: $[-0.0186,\,0.0509]$); the corresponding
micro-$F_1$ intervals also included zero.

Under unseen conditions, the recombination advantage increased. Relative to
one-sided context zero-shot, recall improved by 0.0463
(95\% CI: $[0.0231,\,0.0694]$) and micro-$F_1$ by 0.0413
(95\% CI: $[0.0223,\,0.0636]$). Relative to two-sided zero-shot, the
corresponding improvements were 0.0648
($[0.0278,\,0.1065]$) and 0.0554
($[0.0232,\,0.0895]$), respectively. One-sided--two-sided differences
remained unresolved for both metrics.

At the physical-run level, the recombination-versus-two-sided
difference-of-differences was 0.0400 for micro-$F_1$
(95\% CI: $[0.0088,\,0.0731]$), 0.0509 for constituent recall
($[0.0185,\,0.0880]$), and 0.0833 for $R_{\mathrm{all}}$
($[0.0278,\,0.1481]$).

Because the 108 compound runs represent only nine distinct compositions, the
interaction was also evaluated by composition-level cluster bootstrap. Under
this stricter analysis, the micro-$F_1$ interaction remained positive but
unresolved
($\Delta_{\mathrm{int}}=0.0400$, 95\% CI:
$[-0.0035,\,0.0861]$), whereas constituent recall
($0.0509$, $[0.0093,\,0.0926]$) and $R_{\mathrm{all}}$
($0.0833$, $[0.0093,\,0.1574]$) remained resolved. The micro-$F_1$
interaction was positive in all nine leave-one-composition-out analyses
(0.0236--0.0527), indicating stable direction but composition-level
uncertainty in its magnitude. 

Because recombination retained more training runs than two-sided withholding,
the recombination training sets were repeatedly downsampled to match the
corresponding two-sided training-set cardinalities. Matching was exact in both
condition regimes. Across 30 repeated subsamples, the unseen-condition
recombination advantage remained positive on average for micro-$F_1$
($+0.0335$), constituent recall ($+0.0438$), and
$R_{\mathrm{all}}$ ($+0.0778$). However, composition-level bootstrap
intervals included zero for micro-$F_1$
($[-0.0151,\,0.0813]$), constituent recall
($[-0.0046,\,0.0880]$), and $R_{\mathrm{all}}$
($[-0.0185,\,0.1667]$).

The corresponding context-by-condition interactions also remained positive
on average but unresolved: $+0.0105$ for micro-$F_1$
(95\% CI $[-0.0401,\,0.0606]$), $+0.0205$ for constituent recall
($[-0.0370,\,0.0880]$), and $+0.0312$ for $R_{\mathrm{all}}$
($[-0.0741,\,0.1574]$). Thus, the directional pattern of the primary
analysis was retained after matching total training-set cardinality, but the
interaction was not statistically resolved under this sensitivity analysis.
Accordingly, the primary context effect should not be interpreted as
independent of available training support.

The recombination-versus-one-sided interaction was unresolved at the
run-cluster level for micro-$F_1$ and recall, although $R_{\mathrm{all}}$
showed a positive interaction of 0.0463
(95\% CI: $[0.0046,\,0.1019]$). In the primary analysis, complete
compound-context deprivation had a larger adverse effect under
operating-condition shift. However, after matching total training-set
cardinality, the recombination-versus-two-sided interaction remained
directionally positive but unresolved for all three principal recovery
outcomes.

One-sided directional results were also heterogeneous across the directional
withholding subsets. When broken-bar compound context was withheld, aggregate
constituent recall within the corresponding test subset was 0.979 and 0.896
under seen and unseen conditions, respectively. The corresponding subset-level
values were 0.438 and 0.479 when winding context was withheld and 0.500 under
both condition settings when dynamic-eccentricity context was withheld. These
values are aggregate recall over all true constituents in the affected
compound runs and should not be interpreted as recall of the withheld
constituent itself. In particular, the mechanism-specific analysis in
Table~\ref{tab:constituent_failure_modes} shows that mean
dynamic-eccentricity recall was only 0.017 across the crossed scenarios,
despite near-perfect isolated-fault discrimination.
Both directional predictions from each physical run remained within the same
bootstrap cluster.

\begin{figure*}[t]
\centering
\includegraphics[width=\textwidth]{Fig5_paired_protocol_differences.png}
\caption{Paired run-cluster bootstrap differences between context regimes.
Points denote differences in constituent recall and micro-$F_1$; horizontal
bars denote 95\% bootstrap confidence intervals. Positive values favor the
first regime in each contrast.}
\label{fig:paired_protocol_differences}
\end{figure*}

\subsection{Sensitivity to Decision-Threshold Objective}
\label{subsec:threshold_sensitivity}

Post-hoc threshold calibration with $F_{0.5}$ and balanced accuracy was used
to assess operating-point sensitivity. Relative to the locked $F_1$
objective, $F_{0.5}$ reduced FA/run from 0.639--0.732 to 0.222--0.389 and
increased exact-match recovery from 0.083--0.157 to 0.157--0.287, while
reducing constituent recall to 0.583--0.644. Balanced-accuracy calibration
shifted the trade-off toward higher recall and more false additions. Thus,
exact-set recovery was substantially more operating-point sensitive than
aggregate recall, and the locked $F_1$ operating point should not be viewed
as an upper bound on achievable decomposition performance.

\subsection{Sensitivity to Single-Fault Severity Merging}
\label{subsec:severity_results}

Removing low-severity single-fault support did not produce a uniform change
in constituent recall, exact-match recovery, or all-constituent recovery
across the crossed regimes. Most paired differences for these metrics had
95\% confidence intervals spanning zero, indicating that the principal
generalization findings were not driven by pooling low- and high-severity
single-fault examples.

The clearest differences were observed in false additions. Under seen
conditions, removing low-severity support increased FA/run from 0.653 to
0.755 in the one-sided regime
($\Delta=+0.102$, 95\% CI: $[0.046,\,0.162]$) and from 0.639 to 0.722 under
recombination
($\Delta=+0.083$, 95\% CI: $[0.019,\,0.157]$).
Under unseen conditions, the two-sided regime similarly increased from
0.731 to 0.843 FA/run
($\Delta=+0.111$, 95\% CI: $[0.037,\,0.185]$).

For seen-condition recombination, micro-$F_1$ decreased from 0.677 to 0.651
when low-severity support was removed
($\Delta=-0.026$, 95\% CI: $[-0.052,\,-0.002]$). Other micro-$F_1$
differences were unresolved. Overall, the sensitivity analysis indicates that
severity merging does not account for the principal crossed-generalization
pattern and, in several regimes, the additional low-severity support reduced
over-diagnosis rather than artificially inflating constituent recovery.

\subsection{Robustness Across Classifier Families}
\label{subsec:classifier_robustness_results}

The crossed protocol was repeated with random forest, RBF-SVM, and MLP
classifiers as post-hoc robustness analyses. Table~\ref{tab:model_robustness}
summarizes descriptive performance averaged across the six crossed scenarios;
logistic regression remained the locked primary model.

\begin{table}[t]
\caption{Descriptive classifier-family robustness averaged across the six
crossed evaluation scenarios. LR denotes logistic regression and MLP denotes
the three-hidden-layer feed-forward neural baseline.}
\label{tab:model_robustness}
\centering
\footnotesize
\begin{tabular}{lcccc}
\toprule
Model &
Micro-$F_1$ &
EM &
$R_{\mathrm{const}}$ &
FA/run \\
\midrule
LR      & 0.660 & 0.120 & 0.660 & 0.680 \\
RF      & 0.744 & 0.282 & 0.669 & 0.256 \\
RBF-SVM & 0.692 & 0.157 & 0.664 & 0.509 \\
MLP     & 0.737 & 0.310 & 0.659 & 0.256 \\
\bottomrule
\end{tabular}
\end{table}

Absolute decomposition performance varied substantially by classifier family:
RF achieved the highest mean micro-$F_1$ (0.744), MLP the highest exact-match
rate (0.310), and both reduced mean FA/run to 0.256, while constituent recall
remained similar across models (0.659--0.669). Nevertheless, the
unseen-condition recombination advantage over two-sided context zero-shot
remained positive for all four classifiers
(micro-$F_1$: 0.023--0.055; constituent recall: 0.032--0.069), indicating
that the qualitative context effect was not specific to the locked linear probe.

\subsection{Oracle Detectability and Cross-Partner Constituent Transfer}
\label{subsec:oracle_transfer_results}

Within-pair analysis first tested whether constituent evidence remained
measurable in the target compound. Across all 18 directed relations $A|B$,
condition-held-out AUC for the contrast $A+B$ versus $B$ ranged from 0.9167
to 1.000, with a median of 1.000. Thus, target compounds remained strongly
separable from their partner-only references.

However, recording-related nuisance variation was also detectable. The two
available recordings of the same bearing-outer fault under the same operating
condition were nearly perfectly separable ($\mathrm{AUC}=0.998$), indicating
that high within-pair AUC cannot be attributed uniquely to the added
constituent.

A same-date sensitivity control reduced recording-date mismatch. Eight
same-date compound pairs sharing one constituent and all 12 operating
conditions yielded 16 directed comparisons, with leave-one-condition-out AUCs
of 0.972--1.000. Static eccentricity in the bearing-outer context remained
perfectly separable from same-date alternatives containing either broken bar
or winding ($\mathrm{AUC}=1.000$). These controls do not causally isolate the
target mechanism, but they show that compound-state separability cannot be
explained solely by recording-date mismatch.

Cross-partner transfer therefore tested a stricter question: whether evidence
for constituent $A$ learned from other compound contexts, such as $A+C$
versus $C$, generalized to the target contrast $A+B$ versus $B$. Under seen
conditions, directional cross-partner AUC across the 18 relations had a median
of 0.905 and a mean of 0.815; 13 relations achieved AUC $\geq0.80$ and 15
achieved AUC $\geq0.70$. Under unseen conditions, the median and mean remained
0.884 and 0.803, respectively. Selected results are shown in
Table~\ref{tab:transfer_examples}, and the full directional pattern in
Fig.~\ref{fig:cross_partner_matrix}.

\begin{table*}[t]
\caption{Selected directed cross-partner transfer results. Within-pair AUC
quantifies detectability of constituent $A$ in the target pair, whereas
cross-partner AUC quantifies transfer of evidence learned from other compound
partners.}
\label{tab:transfer_examples}
\centering
\footnotesize
\begin{tabular}{llllccc}
\toprule
Target composition &
Constituent $A$ &
Partner $B$ &
\makecell{Within-pair\\AUC} &
\makecell{Cross-partner\\AUC (seen)} &
\makecell{Cross-partner\\AUC (unseen)} \\
\midrule

Bearing inner + bearing outer
& Bearing outer
& Bearing inner
& 1.000
& 1.000
& 1.000 \\

Bearing inner + dynamic eccentricity
& Bearing inner
& Dynamic eccentricity
& 1.000
& 0.944
& 0.931 \\

Bearing inner + broken bar
& Broken bar
& Bearing inner
& 1.000
& 0.893
& 0.851 \\

Bearing inner + dynamic eccentricity
& Dynamic eccentricity
& Bearing inner
& 1.000
& 0.711
& 0.653 \\

Bearing outer + static eccentricity
& Static eccentricity
& Bearing outer
& 0.917
& 0.000
& 0.000 \\

Bearing outer + winding
& Winding
& Bearing outer
& 0.917
& 0.333
& 0.340 \\

\bottomrule
\end{tabular}
\end{table*}

Transfer was nevertheless strongly heterogeneous. For static eccentricity with bearing outer, within-pair AUC was 0.9167, whereas directional cross-partner AUC was 0.000 under both seen and unseen conditions. Because $\mathrm{AUC}_{\mathrm{sep}}=1.000$, the target remained perfectly rank separable but in the opposite orientation. This was the only relation with a statistically supported score-orientation reversal after multiplicity correction (Holm-adjusted $p=0.0136$ seen; $p=0.0108$ unseen). Given the recording separability observed above, this result is interpreted as reproducible context-dependent score reorientation rather than causal proof of physical signature inversion.

Winding with bearing outer also produced below-chance directional AUCs (0.3333 seen; 0.3403 unseen), but orientation-free separability was only approximately 0.67 and the Holm-adjusted tests were not significant.
Accordingly, this relation was classified as orientation-ambiguous rather than as a supported reversal.

\begin{figure*}[t]
\centering
\includegraphics[width=\textwidth]{Fig2_cross_partner_matrix.png}
\caption{Directional cross-partner transfer AUC. Rows denote constituent $A$
and columns partner $B$. Values below 0.5 indicate below-chance
fixed-orientation ranking but do not alone establish a statistically supported
score-orientation reversal. $n_C$ denotes the number of eligible compound
training contexts for the recovered constituent.}
\label{fig:cross_partner_matrix}
\end{figure*}

Overall, compound states remained strongly discriminable while score
orientation learned from other compound partners could fail to transfer.
Because recording-specific structure was also detectable, the result is
interpreted as context-dependent score transfer rather than causal evidence
of physical constituent-signature inversion. Figure~\ref{fig:within_cross_transfer}
contrasts within-pair detectability with directional cross-partner transfer.

\begin{figure}[H]
\centering
\includegraphics[width=1\columnwidth]{Fig3_within_vs_cross.png}
\vspace{-1mm}
\caption{Within-pair detectability versus directional cross-partner transfer.
The dashed line denotes chance-level ranking. Below-chance values require the
separate score-orientation analysis in
Section~\ref{subsec:statistical_analysis} before being classified as supported
reversals.}
\label{fig:within_cross_transfer}
\end{figure}

\subsection{Physical-Signature Heterogeneity and Directional Failure Modes}
\label{subsec:failure_modes}
Feature-level preservation showed that score reorientation was not equivalent to uniform inversion of the underlying physical descriptors. For static eccentricity with bearing outer, the median preservation ratio was 0.760 over the full ratio-eligible set; among ratio-eligible observations with positive single-fault evidence, 17.8\% showed sign reversal and 34.6\% severe suppression. Conversely, bearing-outer evidence in the bearing-inner compound
had a lower median preservation ratio ($\rho=0.305$); within the corresponding positive-single-evidence subset, 18.4\% of observations showed sign reversal and 58.4\% severe suppression, despite a cross-partner AUC of 1.000. Thus, feature suppression and multivariate score reorientation represent distinct partner-dependent phenomena (Fig.~\ref{fig:signature_preservation}).

\begin{figure}[!t]
\centering
\includegraphics[width=\columnwidth]{Fig4_preservation_polarity.png}
\caption{Feature-level preservation of selected static-eccentricity
signatures under bearing-inner and bearing-outer partner contexts. Features
were selected by prioritizing sign reversals and then partner-dependent
preservation differences; aggregate preservation ratios were calculated over the full ratio-eligible
feature--condition set, whereas sign-reversal and severe-suppression rates were
calculated only over ratio-eligible observations with positive single-fault
evidence. Negative values indicate sign reversal, values
near one preservation, and values above one amplification.}
\label{fig:signature_preservation}
\end{figure}

The locked probe also produced 0.64--0.73 false additions per prediction row,
indicating a material over-diagnosis burden despite partial constituent
recovery. Recording-related nuisance variation further limits physical
attribution: distinct recordings of the same fault state were strongly
separable, while recording date provides only a coarse proxy for experimental
session. Accordingly, the cross-partner results support context-dependent
score transfer but not causal attribution to physical signature inversion.

Overall, the experiments reveal complementary forms of generalization
failure. Mechanism-specific analysis shows that otherwise learnable
constituents can degrade sharply in compound evaluation, with substantially
different patterns across mechanisms. The crossed protocol additionally
suggests that complete compound-context deprivation is more detrimental when
the operating condition is also unseen in the primary analysis, although this
interaction becomes unresolved after matching total training-set cardinality.
Finally, cross-partner analysis shows that constituent scores can remain
strongly separable yet fail to preserve their learned orientation across
partners. These observations describe related but distinct limitations of
constituent generalization rather than a single uniform failure mode.

\section{Conclusion}
\label{sec:conclusion}

This paper evaluated whether the elementary mechanisms of a compound motor
fault can be recovered when their exact combination, their compound context,
and the operating condition are varied independently of the training history.
On a public multimodal induction-motor dataset whose compound faults span
mechanical and electrical modalities, a locked linear constituent probe over
mechanism-specific physical features showed that strong closed-set
discriminability (24-class accuracy 0.913, 95\% CI $[0.885, 0.938]$) coexists
with exact constituent-set recovery of only 0.083--0.157 on unseen compounds.
The two results are not in tension; they measure different abilities, and the
gap between them is the paper's subject.

Four conclusions follow from the controlled analyses. First, at the aggregate level, the dominant failure mode is incomplete
decomposition rather than complete diagnostic silence. Constituent recall remained between
0.630 and 0.694, but a prevalence-only constant predictor could itself reach
0.667 recall at the cost of 1.667 false additions per run and zero exact
matches. Under recombination, all-constituent recovery reached 0.417--0.444, compared with 0.333 for the reported three-constituent prevalence reference; however,
that reference incurred 1.667 false additions per run compared with
approximately 0.64--0.66 for the learned detector. These results show why constituent recall must be
interpreted jointly with false additions, micro-$F_1$, and exact-set recovery
rather than as a stand-alone measure of diagnostic skill. The constituent-resolved controls further showed that these failures do not
share a common origin. Dynamic eccentricity was almost perfectly learnable as
an isolated fault (AUROC 0.9995; $F_1=0.923$) but collapsed to compound
recall of 0.017 and $F_1$ of 0.028, providing the clearest evidence that poor
compound recovery cannot be explained by weak isolated detectability alone.
Static eccentricity was likewise highly learnable in isolation and retained
substantial compound ranking information, while its $F_1$ increased from
0.577 with the locked linear probe to 0.866 with RBF-SVM, consistent with
pronounced classifier-family sensitivity. Winding showed a third pattern: isolated-fault
specificity was already weaker and none of the nonlinear alternatives rescued
its poor compound performance. In contrast, bearing-inner, bearing-outer, and
broken-bar evidence remained strongly discriminative across both isolated and
compound settings. Thus, constituent-level failure in this dataset exhibits distinct empirical
patterns consistent with compound degradation, classifier-capacity
sensitivity, and limited constituent specificity rather than a uniform cross-modal deficit. Second, the adverse effect of compound-context deprivation was larger when
the operating condition was simultaneously unseen in the primary analysis.
For the recombination-versus-two-sided interaction, composition-level
clustering yielded $+0.040$ for micro-$F_1$
(95\% CI $[-0.004,\,0.086]$), $+0.051$ for constituent recall
(95\% CI $[0.009,\,0.093]$), and $+0.083$ for all-constituent recovery
(95\% CI $[0.009,\,0.157]$). Thus, the primary composition-level analysis
supported the interaction for constituent recall and all-constituent recovery,
while the micro-$F_1$ interaction remained unresolved. After matching
recombination training-set cardinality to the two-sided regime, the
interactions remained positive on average but were unresolved for micro-$F_1$
($+0.010$, 95\% CI $[-0.040,\,0.061]$), constituent recall
($+0.021$, $[-0.037,\,0.088]$), and all-constituent recovery
($+0.031$, $[-0.074,\,0.157]$). The observed context effect is therefore
sensitive to available training support and should not be interpreted as a
training-cardinality-independent effect of compound-context exposure. Third,
constituent evidence transfers across compound partners often but
directionally: within-pair detectability was near-perfect throughout, yet
cross-partner transfer ranged from perfect to a statistically supported
score-orientation reversal (Holm-adjusted $p \le 0.014$)---and the
demonstrated separability of distinct recordings of the same fault state
requires that reversal to be interpreted as context-dependent score
transfer rather than as proven physical signature inversion. Fourth,
absolute decomposition quality is an operating-point and model property:
precision-weighted thresholding raised exact-match recovery to 0.157--0.287
at reduced recall, and nonlinear classifier families reduced false additions
from ${\sim}0.68$ to ${\sim}0.26$ per run while the unseen-condition recombination advantage over two-sided withholding
remained directionally positive across all four classifier families.

For practice, these results caution against reading closed-set accuracy as
deployment readiness, argue for calibrating constituent detectors with
explicit precision--recall cost structure, and indicate that commissioning
data should sample compound contexts and operating conditions jointly rather
than relying on either alone. Several directions remain open. The dataset provides insufficient replicated same-state recordings to
separate fault, assembly, sensor-mounting, thermal, and recording-related
effects. The transfer analyses therefore stop at carefully bounded
attribution; systematically replicated acquisitions would enable stronger
separation of these potential sources of variation. Representation or normalization strategies that promote
constituent-score consistency across compound partners are identified as
hypotheses, not established remedies. Finally, the constituent-level
protocol, prevalence references, and controls used here are directly
portable to the companion gearbox release and to other multimodal datasets,
where cross-dataset replication of the context-by-condition interaction
would test its generality.

\section*{Data Availability}
The dataset analyzed in this study is publicly available from Mendeley Data
at \url{https://doi.org/10.17632/6s3dggj9mw.1} under a CC~BY~4.0 license
\cite{mcc5motor_data}, with its accompanying descriptor article
\cite{chen2026}. The acquisition-timestamp metadata used for the
recording-sensitivity and same-date controls were recovered from the released
file names and are distributed with the analysis code.

\section*{Code Availability}
The analysis code, the derived run-level feature tables, the split
definitions for every protocol reported here, and the scripts that regenerate
each table and figure are publicly available at
\url{https://doi.org/10.5281/zenodo.XXXXXXX}.

\bibliographystyle{IEEEtran}
\bibliography{references}

\begin{IEEEbiography}[{\includegraphics[width=1in,height=1.25in,clip,keepaspectratio]{evin.png}}]{Evin \c{S}ahin Sadik} received the B.Sc. degree in Electrical and Electronics Engineering from Bilecik \c{S}eyh Edebali University, Bilecik, Turkey, in 2012, and the M.Sc. and Ph.D. degrees in Electrical and Electronics Engineering from K\"utahya Dumlup{\i}nar University, K\"utahya, Turkey, in 2016 and 2022, respectively. She is currently an Assistant Professor with the Department of Electrical and Electronics Engineering, K\"utahya Dumlup{\i}nar University. Her research interests include signal processing, machine learning, and deep learning.
\end{IEEEbiography}

\begin{IEEEbiography}[{\includegraphics[width=1in,height=1.25in,clip,keepaspectratio]{Huseyin_Canseven.png}}]{H{\"u}seyin Tayyer Canseven} (Member, IEEE) received the B.Sc. and M.Sc. degrees in Electrical and Electronics Engineering from Ege University, Izmir, Turkey, in 2018 and 2022, respectively, and the D.Sc. degree in Electrical Engineering from Lappeenranta-Lahti University of Technology (LUT), Lappeenranta, Finland, in 2025. He is currently a Post-Doctoral Researcher with the Department of Electrical Engineering, LUT University. His research interests focus on design and condition monitoring of electrical machines and motor drives.
\end{IEEEbiography}

\EOD

\end{document}